\section{The Moment No One Occupies}

Augustine, having defined death as a parting, went looking for the instant at which the parting occurs, and could not find it. The chapters in which he fails are among the strangest in the \emph{City of God}, and they are not a digression.

The problem is stated with almost mathematical dryness. There are three times---before death, in death, after death---and three corresponding states: living, dying, dead. Now take a man in extremity. ``So long as the soul is in the body, especially if consciousness remain, the man certainly lives; for body and soul constitute the man. And thus, before death, he cannot be said to be in death; but when, on the other hand, the soul has departed, and all bodily sensation is extinct, death is past, and the man is dead. Between these two states the dying condition finds no place; for if a man yet lives, death has not arrived; if he has ceased to live, death is past. Never, then, is he dying.''\endnote{Augustine, \emph{De civitate Dei} XIII.11, Dods 1871, chapter heading ``Whether one can both be living and dead at the same time''; verified 2026-07-26 in the source library's tracked transcription of the 1871 printing, volume~1. The parallel arguments are at XIII.9 (``Whether we should say that the moment of death, in which sensation ceases, occurs in the experience of the dying or in that of the dead'') and XIII.10.} And he presses it further, into the place where it becomes uncomfortable: ``it is difficult to explain how we speak of those who are not yet dead, but are agonized in their last and mortal extremity, as being in the article of death. Yet what else can we call them than dying persons? \ldots{} No one, therefore, is dying unless living.''\endnote{\emph{De civitate Dei} XIII.9, Dods 1871; verified as above. Augustine's conclusion in the same chapter: ``The same person is therefore at once dying and living, but drawing near to death, departing from life; yet in life, because his spirit yet abides in the body; not yet in death, because not yet has his spirit forsaken the body.''}

The comparison he reaches for is the present tense: ``you try to lay your finger on the present, and cannot find it, because the present occupies no space, but is only the transition of time from the future to the past.'' The instant of death is like that---a boundary rather than a duration, real as a limit and unavailable as an experience. Which produces the conclusion he states and immediately declines to accept: ``Must we then conclude that there is thus no death of the body at all? For if there is, where is it, since it is in no one, and no one can be in it?''\endnote{\emph{De civitate Dei} XIII.11, Dods 1871; verified as above. Augustine's resolution is linguistic rather than metaphysical: ``Let us, then, speak in the customary way,---no man ought to speak otherwise,---and let us call the time before death come, `before death'\ldots{} And let us use the same phraseology as Scripture uses; for it makes no scruple of saying that the dead are not after but in death,'' citing Psalm 6:6 (``For in death there is no remembrance of thee''). The Douay renders that verse ``For there is no one in death, that is mindful of thee: and who shall confess to thee in hell?'' (read 2026-07-26 in the tracked Challoner verse text), which reads the Latin \latin{in morte} as a state one is in---exactly Augustine's point about scriptural usage.}

He does not resolve it. He decides to keep speaking the way everyone speaks, and he gives a reason that is better than a resolution: Scripture itself talks about being \emph{in} death, so the ordinary usage has warrant even where the analysis fails. This is Augustine at his best---an argument he cannot answer, reported as an argument he cannot answer, followed by a working practice.

Between the two chapters where he presses the puzzle he inserts one that has been quoted more often than either, and it is the reason the puzzle is not merely a curiosity. ``For no sooner do we begin to live in this dying body, than we begin to move ceaselessly towards death\ldots{} our whole life is nothing but a race towards death, in which no one is allowed to stand still for a little space, or to go somewhat more slowly, but all are driven forwards with an impartial movement, and with equal rapidity. For he whose life is short spends a day no more swiftly than he whose life is longer.''\endnote{\emph{De civitate Dei} XIII.10, Dods 1871, chapter heading ``Of the life of mortals, which is rather to be called death than life''; verified as above. Augustine's arithmetic is exact and worth noticing: the dying do not travel faster, they have less ground to cover---``It is one thing to make a longer journey, and another to walk more slowly.''} The instant is unlocatable because the process is not instantaneous. If dying is something a body has been doing since it began, then asking when it \emph{happens} is like asking at what moment a river arrives at the sea.

\subsection{What the Church does about an instant it cannot find}

The interesting thing is what the Church, which has never resolved Augustine's puzzle, does about it in practice. It legislates.

Canon 1004 of the current Latin code provides that the anointing of the sick may be administered to a member of the faithful ``who, having reached the use of reason, begins to be in danger due to sickness or old age.'' Then canon 1005: ``This sacrament is to be administered in a case of doubt whether the sick person has attained the use of reason, is dangerously ill, or is dead.''\endnote{\emph{Codex Iuris Canonici} (1983), cann.~1004 §1 and 1005; English from the Holy See's web text of the Code, read 2026-07-26. The governing body of law is the 1983 Code of Canon Law for the Latin Church, in force since 27 November 1983; the parallel Eastern discipline is in the \emph{Codex Canonum Ecclesiarum Orientalium} and is not treated here. The canon is disciplinary law, not doctrine, and states a rule of practice, not a theory of death.}

Read the third clause slowly. The Church has legislated for the case in which the minister does not know whether the person in front of him is alive. It does not tell him to determine the fact; it tells him to act. The sacrament goes forward under doubt, conditionally, on the side of the person.

That is not an accident of drafting, and it is not merely charitable improvisation. It is what a body has to do when it holds two things at once: that the sacraments are for the living, and that the boundary of living is genuinely not locatable from outside. If the moment were findable, the canon would be superfluous and slightly scandalous---one does not anoint corpses. If the moment were unreal, the canon would be incoherent. The canon exists because the moment is real and unfindable, which is precisely Augustine's result.

\begin{studybox}{Project synthesis: law at the boundary Augustine could not draw}
The claim here is this essay's own, argued from the texts cited and attributed to no source: the Church's disciplinary practice at the deathbed is best read as a jurisprudence built on top of an admitted metaphysical indeterminacy, and the two halves are more coherent together than either is alone. Augustine established that no observer, and not even the dying man, occupies the instant of death; the canonical rule instructs the minister to act \emph{as if} the person may still be within reach whenever he cannot know otherwise. The relation between them is not deduction---no canonist cites \emph{De civitate Dei} XIII.11 in the margin of canon 1005---but the fit is exact, and it explains a feature of Catholic deathbed practice that puzzles observers: the priest who anoints a body that has apparently stopped is not confused about biology. He is declining to place a bet at the one point in a human life where the tradition has always said the bet cannot be placed.

\medskip\noindent\textbf{The strongest objection.} The canon is simply prudential and requires no such reading. Sacramental law is full of provisions for doubt---doubtful baptism, doubtful matter, doubtful consent---and they express nothing more than the ordinary principle that in doubt the safer course favors the recipient's salvation. On this account canon 1005 is a routine application of \latin{in dubio pro salute animae} and carries no metaphysical commitment about the instant of death whatever; reading Augustine into it is decoration. The objection is strong because it is the simpler explanation and because the canon really does sit among other doubt-provisions that have nothing to do with the boundary of life.

\medskip\noindent\textbf{What would change the judgment.} If the Church's own commentary on this canon or its predecessors grounded the doubt-clause solely in the uncertainty of \emph{medical observation}---that is, in our not yet knowing how to tell, with the implication that better instruments would remove the doubt---the synthesis here would be wrong, and the canon would be a temporary accommodation to ignorance rather than a standing accommodation to a boundary. The essay has not examined the drafting history of canon 1005 or of its predecessor in the 1917 Code, which it did not consult, and records this as an open question rather than a settled point.
\end{studybox}

\subsection{Where this essay stops}

One boundary must be marked here rather than later, because it is the place where readers most often expect a theological essay to overreach. Nothing said above settles any contemporary question about the determination of death---whether neurological criteria are adequate, when organs may be removed, what a family should conclude from a particular prognosis. Those are questions of fact and of applied moral theology, they turn on medical evidence this essay has not examined, and the Church's interventions in them belong to a different order of teaching than the definitions quoted above.\endnote{The essay takes no position on neurological or circulatory criteria for the determination of death, on organ procurement, or on the withdrawal of treatment. Those questions require medical evidence and a moral-theological treatment outside the scope of this study; a reader facing them should consult the competent pastoral and medical authorities with the complete facts. The definitional teaching quoted above---that death is the separation of soul and body---does not by itself yield a criterion, and the Church has never claimed that it does.} What follows from Augustine and from canon 1005 is narrower and more durable: the philosophical definition of death does not hand anyone a test for it, and the Church knows this, and has arranged its practice accordingly.

The definition, then, is secure and the instant is not. That is an odd result, and it is about to be joined by a second and larger oddity, this time not about \emph{when} death happens but about \emph{why}.
